\documentclass[11pt,a4paper]{article}

% ============================================================
% Packages
% ============================================================

\usepackage[utf8]{inputenc}
\usepackage[T1]{fontenc}
\usepackage{lmodern}

\usepackage{amsmath}
\usepackage{amssymb}
\usepackage{amsthm}
\usepackage{mathtools}
\usepackage{bm}

\usepackage{geometry}
\geometry{
  left=30mm,
  right=30mm,
  top=30mm,
  bottom=30mm
}

\usepackage{hyperref}
\usepackage{enumitem}

% ============================================================
% Theorem environments
% ============================================================

\newtheorem{definition}{Definition}[section]
\newtheorem{proposition}[definition]{Proposition}
\newtheorem{theorem}[definition]{Theorem}
\newtheorem{remark}[definition]{Remark}

% ============================================================
% Commands
% ============================================================

\newcommand{\R}{\mathbb{R}}
\newcommand{\C}{\mathbb{C}}
\newcommand{\Z}{\mathbb{Z}}
\newcommand{\ii}{\mathrm{i}}
\newcommand{\dd}{\mathrm{d}}
\newcommand{\Om}{\Omega}
\newcommand{\Gam}{\Gamma}
\newcommand{\End}{\operatorname{End}}
\newcommand{\id}{\operatorname{id}}
\newcommand{\vol}{\operatorname{vol}}
\newcommand{\PD}{\operatorname{PD}}

% ============================================================
% Title
% ============================================================

\title{\textbf{Maxwell Theory from a Hermitian Complex Line Bundle}}
\author{}
\date{}

% ============================================================
% Document
% ============================================================

\begin{document}

\maketitle
\tableofcontents

\section{Spacetime}

We begin with a smooth, oriented, four-dimensional manifold
\[
M.
\]

To formulate classical electromagnetism, we equip $M$ with a
Lorentzian metric
\[
g \in \Gamma\bigl(S^2T^*M\bigr)
\]
of signature
\[
(-,+,+,+).
\]

Thus, at each point $x\in M$, the metric defines a nondegenerate
symmetric bilinear form
\[
g_x:T_xM\times T_xM\longrightarrow \R.
\]

We assume that $M$ is oriented and time-oriented.

The metric and orientation determine the Hodge star operator
\[
*:\Omega^k(M)\longrightarrow\Omega^{4-k}(M).
\]

For differential forms $\alpha,\beta\in\Omega^k(M)$, the Hodge star
is characterized by
\[
\alpha\wedge *\beta
=
\langle \alpha,\beta\rangle_g\,\vol_g,
\]
where $\vol_g$ denotes the metric volume form.

The Hodge star is the part of Maxwell theory that depends explicitly
on the geometry of spacetime.

\section{Complex Line Bundles}

\begin{definition}[Complex line bundle]
A \emph{complex line bundle} over $M$ is a smooth vector bundle
\[
\pi:L\longrightarrow M
\]
whose fibers
\[
L_x=\pi^{-1}(x)
\]
are one-dimensional complex vector spaces.

Equivalently,
\[
\dim_{\C}L_x=1
\]
for every $x\in M$.
\end{definition}

Locally, every complex line bundle is trivial. Thus for every
$x\in M$, there exists an open neighborhood $U\subset M$ such that
\[
L|_U\cong U\times\C.
\]

A smooth section of $L$ is a smooth map
\[
\psi:M\longrightarrow L
\]
such that
\[
\pi\circ\psi=\id_M.
\]

The space of smooth sections is denoted by
\[
\Gamma(L).
\]

\section{Hermitian Structure}

\begin{definition}[Hermitian structure]
A \emph{Hermitian structure} on $L$ is a smoothly varying family of
positive-definite Hermitian inner products
\[
h_x:L_x\times L_x\longrightarrow\C.
\]

Thus
\[
h_x(v,w)=\overline{h_x(w,v)}
\]
and
\[
h_x(v,v)>0
\]
for every nonzero $v\in L_x$.
\end{definition}

A complex line bundle equipped with such a metric is called a
\emph{Hermitian complex line bundle}.

The Hermitian structure reduces the structure group of the line
bundle from
\[
\C^\times
\]
to
\[
U(1).
\]

Consequently, electromagnetism may be interpreted geometrically as
a $U(1)$ gauge theory.

\section{Connections}

Values of a section at distinct spacetime points belong to different
vector spaces:
\[
\psi(x)\in L_x,
\qquad
\psi(y)\in L_y.
\]

There is no canonical way to compare these vectors. A connection
provides an infinitesimal rule for doing so.

\begin{definition}[Connection]
A connection on $L$ is a $\C$-linear map
\[
\nabla:\Gamma(L)
\longrightarrow
\Omega^1(M;L)
\]
satisfying the Leibniz rule
\[
\nabla(f\psi)
=
\dd f\otimes\psi
+
f\nabla\psi
\]
for every
\[
f\in C^\infty(M,\C),
\qquad
\psi\in\Gamma(L).
\]
\end{definition}

\begin{definition}[Hermitian connection]
A connection $\nabla$ is compatible with the Hermitian structure
if
\[
\dd\,h(\psi,\phi)
=
h(\nabla\psi,\phi)
+
h(\psi,\nabla\phi)
\]
for all sections $\psi,\phi\in\Gamma(L)$.
\end{definition}

Choose a local unitary frame $e$ over an open subset $U\subset M$.
Every section can then be written locally as
\[
\psi=\psi_U e.
\]

A Hermitian connection takes the local form
\[
\nabla
=
\dd+\ii A,
\]
where
\[
A\in\Omega^1(U;\R).
\]

The real one-form $A$ is the electromagnetic gauge potential.

\section{Gauge Transformations}

A different unitary frame is related to the original one by
\[
e'
=
e^{-\ii\lambda}e,
\]
where
\[
\lambda:U\longrightarrow\R.
\]

Under this change of frame, the local gauge potential transforms as
\[
A
\longmapsto
A+\dd\lambda.
\]

This is precisely the usual $U(1)$ gauge transformation.

Therefore, the potential $A$ is not itself an invariant geometric
object. Rather, it is a local representative of a connection on
the line bundle.

\section{Curvature}

\begin{definition}[Curvature]
The curvature of a connection $\nabla$ is the two-form with values
in $\End(L)$ defined by
\[
F_\nabla
=
\nabla^2.
\]
\end{definition}

For a Hermitian line bundle, the curvature takes values in the Lie
algebra
\[
\mathfrak{u}(1)=\ii\R.
\]

With the convention
\[
\nabla=\dd+\ii A,
\]
we obtain
\[
F_\nabla
=
\ii\,\dd A.
\]

It is convenient to define the real electromagnetic field-strength
two-form by
\[
F=\dd A.
\]

Thus
\[
F_\nabla=\ii F.
\]

Under a gauge transformation,
\[
A\mapsto A+\dd\lambda,
\]
we have
\[
F
\mapsto
\dd(A+\dd\lambda)
=
\dd A
=
F.
\]

Hence $F$ is gauge invariant.

\section{The First Maxwell Equation as a Bianchi Identity}

Since locally
\[
F=\dd A,
\]
we immediately obtain
\[
\dd F
=
\dd^2A
=
0.
\]

More intrinsically, this is the Abelian form of the Bianchi identity.

Thus the first half of Maxwell's equations is not a dynamical
equation. It follows from the geometric fact that the electromagnetic
field is the curvature of a $U(1)$ connection:
\[
\boxed{\dd F=0.}
\]

\section{The Maxwell Action}

The dynamics of the connection are determined by an action
functional.

For the free electromagnetic field, define
\[
S[A]
=
-\frac12
\int_M
F\wedge *F,
\]
where
\[
F=\dd A.
\]

Equivalently, in local coordinates,
\[
S[A]
=
-\frac14
\int_M
F_{\mu\nu}F^{\mu\nu}
\sqrt{|\det g|}\,
\dd^4x.
\]

Here
\[
F
=
\frac12
F_{\mu\nu}
\dd x^\mu\wedge\dd x^\nu,
\]
with
\[
F_{\mu\nu}
=
\partial_\mu A_\nu
-
\partial_\nu A_\mu.
\]

\section{Variation of the Maxwell Action}

Consider a variation
\[
A\longmapsto A+\varepsilon a,
\]
where
\[
a\in\Omega^1(M)
\]
has compact support.

Then
\[
F
\longmapsto
F+\varepsilon\,\dd a,
\]
so
\[
\delta F=\dd a.
\]

The variation of the action is
\begin{align*}
\delta S
&=
-\frac12
\delta
\int_M F\wedge *F
\\
&=
-\int_M
\dd a\wedge *F.
\end{align*}

Using
\[
\dd(a\wedge *F)
=
\dd a\wedge *F
-
a\wedge\dd *F,
\]
we obtain
\[
\dd a\wedge *F
=
\dd(a\wedge *F)
+
a\wedge\dd *F.
\]

Therefore,
\[
\delta S
=
-\int_M\dd(a\wedge *F)
-
\int_M a\wedge\dd *F.
\]

Assuming either that $M$ has no boundary or that the variation
vanishes sufficiently rapidly at the boundary, the first term
vanishes by Stokes' theorem. Hence
\[
\delta S
=
-\int_M
a\wedge\dd *F.
\]

Since $a$ is arbitrary, the Euler--Lagrange equation is
\[
\boxed{\dd *F=0.}
\]

\section{Vacuum Maxwell Equations}

Combining the geometric identity
\[
\dd F=0
\]
with the Euler--Lagrange equation
\[
\dd *F=0,
\]
we obtain the vacuum Maxwell equations:
\[
\boxed{
\begin{aligned}
\dd F &=0,\\
\dd *F&=0.
\end{aligned}
}
\]

The two equations have conceptually different origins:
\[
\begin{array}{ccc}
\text{connection}
&\longrightarrow&
\text{curvature}
\\[4pt]
\nabla
&\longrightarrow&
F
\\[4pt]
&&\downarrow
\\[-2pt]
&&
\dd F=0,
\end{array}
\]
whereas
\[
\begin{array}{ccc}
g
&\longrightarrow&
*
\\[4pt]
&&\downarrow
\\[-2pt]
F
&\longrightarrow&
S[F]
\\[4pt]
&&\downarrow
\\[-2pt]
&&
\dd *F=0.
\end{array}
\]

Thus
\[
\dd F=0
\]
is a geometric identity, while
\[
\dd *F=0
\]
is a dynamical field equation.

\section{Electric and Magnetic Fields}

In Minkowski spacetime with coordinates
\[
(x^0,x^1,x^2,x^3)=(t,x,y,z),
\]
the electromagnetic two-form can be decomposed as
\[
F
=
E_i\,\dd x^i\wedge\dd t
+
\frac12
\epsilon_{ijk}B^k
\,\dd x^i\wedge\dd x^j.
\]

Then
\[
\dd F=0
\]
is equivalent to
\[
\nabla\cdot\mathbf{B}=0
\]
and
\[
\nabla\times\mathbf{E}
+
\frac{\partial\mathbf{B}}{\partial t}
=
0.
\]

Similarly,
\[
\dd *F=0
\]
is equivalent to
\[
\nabla\cdot\mathbf{E}=0
\]
and
\[
\nabla\times\mathbf{B}
-
\frac{\partial\mathbf{E}}{\partial t}
=
0
\]
in units where
\[
c=\epsilon_0=\mu_0=1.
\]

\section{Coupling to an Electric Current}

Let
\[
J\in\Omega^1(M)
\]
denote the electric four-current one-form.

The interaction term may be written as
\[
S_{\mathrm{int}}
=
\int_M A\wedge *J.
\]

The total action is therefore
\[
S[A]
=
-\frac12\int_M F\wedge *F
+
\int_M A\wedge *J.
\]

Variation with respect to $A$ gives
\[
\boxed{
\dd *F=*J.
}
\]

Together with the Bianchi identity,
\[
\boxed{
\dd F=0,
}
\]
the Maxwell equations with sources are
\[
\boxed{
\begin{aligned}
\dd F&=0,\\
\dd *F&=*J.
\end{aligned}
}
\]

Applying $\dd$ to the second equation yields
\[
\dd *J=0.
\]

This is the differential-form expression of electric charge
conservation.

\section{Topology of the Electromagnetic Field}

If the line bundle $L$ is globally trivial, one may choose a global
potential
\[
A\in\Omega^1(M)
\]
and write globally
\[
F=\dd A.
\]

For a nontrivial line bundle, however, a global potential need not
exist.

Instead, one has local potentials
\[
A_\alpha\in\Omega^1(U_\alpha)
\]
related on overlaps by gauge transformations
\[
A_\beta
=
A_\alpha+\dd\lambda_{\alpha\beta}.
\]

Nevertheless, the curvature forms agree:
\[
\dd A_\alpha
=
\dd A_\beta.
\]

Consequently they define a global two-form
\[
F\in\Omega^2(M).
\]

With the appropriate normalization, its de Rham cohomology class
satisfies
\[
\left[
\frac{F}{2\pi}
\right]
=
c_1(L)_{\R},
\]
where
\[
c_1(L)\in H^2(M;\Z)
\]
is the first Chern class of the line bundle.

Thus the topology of the electromagnetic gauge bundle is encoded
in the cohomology class of the electromagnetic field strength.

\section{Geometric Summary}

The construction may be summarized as
\[
\boxed{
(M,g)
\quad\longrightarrow\quad
(L,h)
\quad\longrightarrow\quad
\nabla
\quad\longrightarrow\quad
F_\nabla
\quad\longrightarrow\quad
S[\nabla].
}
\]

More explicitly,
\[
\begin{aligned}
M
&=
\text{spacetime},
\\
L\to M
&=
\text{Hermitian complex line bundle},
\\
\nabla
&=
\text{$U(1)$ connection},
\\
A
&=
\text{local connection one-form},
\\
F
&=
\text{curvature / electromagnetic field},
\\
S
&=
-\frac12\int_M F\wedge *F.
\end{aligned}
\]

The Maxwell equations then arise from two different geometric
principles:
\[
\boxed{
\underbrace{\dd F=0}_{\text{Bianchi identity}},
\qquad
\underbrace{\dd *F=*J}_{\text{Euler--Lagrange equation}}.
}
\]

In this formulation, classical electromagnetism is therefore the
gauge theory of connections on Hermitian complex line bundles.

\section{Relation to Yang--Mills Theory}

The preceding construction is the Abelian case of Yang--Mills
theory.

For electromagnetism the structure group is
\[
U(1),
\]
which is Abelian, and locally
\[
F=\dd A.
\]

For a non-Abelian Lie group $G$, a connection is locally represented
by a Lie-algebra-valued one-form
\[
A\in\Omega^1(M;\mathfrak{g}),
\]
and its curvature becomes
\[
F_A
=
\dd A+A\wedge A.
\]

The Yang--Mills action is
\[
S_{\mathrm{YM}}[A]
=
-\frac12
\int_M
\operatorname{tr}
\left(
F_A\wedge *F_A
\right).
\]

Thus Maxwell theory is the simplest example of the general
geometric principle
\[
\boxed{
\text{gauge field}
=
\text{connection},
\qquad
\text{field strength}
=
\text{curvature}.
}
\]

\end{document}

